% Pursuing Stacks (cotes 134-2 to 134-8): archive cote 134-5, site page 26 (PDF page 27), typescript page 213.
% typescript. Draft compared: ps4.tex:1171-1239.
% AI-transcribed by Claude Opus 5.5 (claude-opus-5-5, adaptive thinking, effort high), Message Batches API, batch ps2 (msgbatch_01NVwr1Yuo8qJf5rykeGBXPK), 2026-11-12; not yet checked by eye.
\scanpage{134-5/26}
\mspage{213}

10.6.

72. In the long last, We'll come back now to the ``asphericity game''\hand{!}
\deleted{again !} Let's take up the exposition at the point where we stopped
two months ago (section 67, p.188 ). We were then about to reformulate
in various ways the property of asphericity, more specifically
\underline{W}-asphericity, for a given map
\[ i : A \longrightarrow B \]
in (Cat). To this end, we introduced the corresponding diagram of maps
in (Cat) with ``commutation morphism'' $\hand{\lambda}_i$ :
\diagram{\begin{tikzcd}[column sep=small]
B^{\wedge} \arrow[rr, "i^{*}"] \arrow[dr, "i_B"', ""{name=iB, right}] & & A^{\wedge} \arrow[dl, "i_A", ""{name=iA, left}] \\
& (\mathrm{Cat}) \arrow[bend right=10, from=iA, to=iB, "\text{\hand{$\lambda_i$}}"'] &
\end{tikzcd}}
\hfill .

As already stated, ,the asphericity condition on $i$ can be expressed \xout{\illegible{1}}
in a variety of equivalent ways, as a condition on either of the three
``aspects''
\[ \hand{\lambda}_i \ , \ i^{*} \ , \ i_A i^{*} \]
of the situation created by $i$, with respect to the localizing sets
$\underline{W}$ , $\underline{W}_A$ , $\underline{W}_B$ in the three categories under consideration, or with res-%
pect to the notion of \underline{W}-aspheric objects in these. As $\underline{W}_A$\hand{,} and aspheric
objects in $A^{\wedge}$\hand{,}are defined \inserted{respectively} in terms of \underline{W} and aspheric objects in (Cat)
via the functor $i_A$ , it turns out that the \xout{characterizations}\inserted[typed]{formulations} in terms of
properties of $i_A i^{*}$ reduce trivially to the corresponding formulations
in terms of $i^{*}$ , which will be the most directly useful for\xout{m} our \xout{pur}
\xout{later} purpose for later work - hence we'll omit them, and focus atten-%
tion instead on $\hand{\lambda}_i$ and $i^{*}$ . Notations are those of loc. cit. , in par-%
ticular we are working with the categories
\[ \mathrm{Hot}^{\underline{W}}_{A} \quad \text{or} \quad \mathrm{Hot}_A \overset{\mathrm{def}}{=} \underline{W}_A^{-1} A^{\wedge} \]
and
\[ \mathrm{Hot}(\underline{W}) = \underline{W}^{-1} (\mathrm{Cat}) \quad , \]
where \underline{W} is a given ``basic localizer'' , namely a set of arrows in (Cat)
satisfying certain conditions.

It seems worthwhile here to be careful to state which exactly are
the properties of \underline{W} we are going to use in the ``asphericity game'' - there%
fore I'm going to list them again here, \xout{\unsure{Once again I under}} using a \deleted{desig-}
\deleted{nation}\inserted{labelling} which\hand{,} hopefully\hand{,} will not have to be changed again:

Loc 1 ) ``Mild saturation'' (cf. page \unsure{\hand{59}|59}) ,

Loc 2 ) ``Homotopy condition'' :$\hand{\Delta}_1 x X \rightarrow X$ is in \underline{W} for any $X$ in (Cat) ,

\margin{26}
